\documentclass[11pt]{article}

\usepackage[margin=1in]{geometry}
\usepackage[T1]{fontenc}
\usepackage[utf8]{inputenc}
\usepackage{amsmath}
\usepackage{amssymb}
\usepackage{enumitem}
\usepackage{verbatim}
\usepackage{hyperref}

\title{Documentation for \texttt{RigidPlaneBoundaryCondition}}
\date{}

\begin{document}

\maketitle

\section{Purpose and Scope}

The \texttt{RigidPlanePenalty} boundary condition implements a unilateral rigid-plane
contact model for displacement-based solid problems. In the current code, it is
supported only for \texttt{struct} and \texttt{ustruct} equations, and it is
assembled as a Neumann-type boundary contribution. The class
\texttt{RigidPlaneBoundaryCondition.cpp} itself stores and validates the rigid-plane
data, while the actual weak-form contribution is assembled elsewhere in the solver.

\section{User Input Parameters}

A boundary is recognized as a rigid-plane penalty boundary when the
boundary-condition type is one of
\begin{verbatim}
RigidPlanePenalty
RigidPlane
Rigid_Plane_Penalty
\end{verbatim}
as parsed in \texttt{Code/Source/solver/read\_files.cpp}.

The following user parameters are required:

\begin{itemize}[leftmargin=*]
\item \texttt{<Type>}:
must be \texttt{RigidPlanePenalty}, \texttt{RigidPlane}, or
\texttt{Rigid\_Plane\_Penalty}.

\item \texttt{<Time\_dependence>}:
must be \texttt{Unsteady}.

\item \texttt{<Temporal\_values\_file\_path>}:
required. The implementation expects a scalar temporal history, which is
interpreted as the rigid-plane displacement amplitude in time.

\item \texttt{<Penalty\_factor>}:
required, must be strictly positive.

\item \texttt{<Plane\_point>}:
required vector of size \texttt{nsd}. This is one point on the reference rigid
plane.

\item \texttt{<Plane\_normal>}:
required vector of size \texttt{nsd}. This defines the plane normal and is
normalized internally by the constructor in
\texttt{RigidPlaneBoundaryCondition.cpp}.
\end{itemize}

The following restrictions are enforced in \texttt{read\_files.cpp}:

\begin{itemize}[leftmargin=*]
\item The boundary condition is allowed only for \texttt{struct} and
\texttt{ustruct}.
\item It cannot be mixed with Robin-specific parameters such as
\texttt{<Spatial\_values\_file\_path>}, \texttt{<Stiffness>},
\texttt{<Damping>}, or \texttt{<Apply\_along\_normal\_direction>}.
\item The temporal history must be scalar.
\item \texttt{Plane\_point} and \texttt{Plane\_normal} must each have exactly
\texttt{nsd} entries.
\end{itemize}

A representative input block is

\begin{verbatim}
<Add_BC>
  <Type> RigidPlanePenalty </Type>
  <Time_dependence> Unsteady </Time_dependence>
  <Temporal_values_file_path> rigid_plane_motion.dat </Temporal_values_file_path>
  <Penalty_factor> 1.0e6 </Penalty_factor>
  <Plane_point> 0.0 0.0 0.0 </Plane_point>
  <Plane_normal> 0.0 1.0 0.0 </Plane_normal>
</Add_BC>
\end{verbatim}

Here, \texttt{rigid\_plane\_motion.dat} provides the scalar plane motion in time,
and the plane is defined by
\[
\mathbf{x}_p = \texttt{Plane\_point},
\qquad
\mathbf{n} = \frac{\texttt{Plane\_normal}}{\|\texttt{Plane\_normal}\|}.
\]

\section{What \texttt{RigidPlaneBoundaryCondition.cpp} Does}

The file \texttt{Code/Source/solver/RigidPlaneBoundaryCondition.cpp} is
responsible for

\begin{itemize}[leftmargin=*]
\item validating the penalty factor,
\item checking that \texttt{Plane\_point} and \texttt{Plane\_normal} are
non-empty and dimensionally consistent,
\item normalizing the plane normal,
\item storing the rigid-plane data,
\item broadcasting the data to all ranks through \texttt{distribute()}.
\end{itemize}

It does not assemble the residual or tangent contributions itself.

\section{Where the Variational Form Is Modified}

The actual weak-form modification is performed in
\[
\texttt{Code/Source/solver/set\_bc.cpp}
\]
inside the function
\[
\texttt{set\_bc\_rigid\_plane\_l(...)}.
\]

The call chain is:

\begin{enumerate}[leftmargin=*]
\item The boundary condition is parsed in \texttt{read\_files.cpp}.
\item During boundary assembly, \texttt{set\_bc\_neu\_l(...)} checks whether
\texttt{lBc.rigid\_plane\_bc.is\_initialized()} is true.
\item If true, it dispatches to \texttt{set\_bc\_rigid\_plane\_l(...)} and
returns.
\end{enumerate}

Therefore, the rigid-plane contact term is assembled in place of the usual
Neumann traction treatment for that boundary face.

\section{Kinematic Definition Used in the Code}

At each boundary quadrature point, the current position is computed as
\[
\mathbf{y} = \sum_{a=1}^{n_{en}} N_a \left(\mathbf{x}_a + \mathbf{d}_a\right),
\]
where \(\mathbf{x}_a\) are nodal reference coordinates and \(\mathbf{d}_a\) are
the interpolated displacements taken from the intermediate displacement state.

The scalar plane motion is read from the temporal function and denoted by
\[
u_p(t).
\]

The code then defines the contact gap as
\[
g = \mathbf{n}\cdot(\mathbf{y}-\mathbf{x}_p) - u_p(t).
\]

This is the quantity implemented in \texttt{set\_bc.cpp} through
\begin{verbatim}
gap = -plane_displacement + plane_normal · (y - plane_point)
\end{verbatim}

The contact condition is unilateral:
\[
g \le 0 \quad \Longrightarrow \quad \text{no contribution},
\]
\[
g > 0 \quad \Longrightarrow \quad \text{penalty traction is applied}.
\]

\section{Penalty Traction and Weak Form}

When the constraint is active (\(g>0\)), the code defines the traction vector
\[
\mathbf{h} = k_p\, g\, \mathbf{n},
\]
where \(k_p = \texttt{Penalty\_factor}\).

The corresponding boundary residual contribution is assembled as
\[
\mathbf{R}_{\Gamma_c}
= \int_{\Gamma_c^{\text{active}}} \mathbf{N}^T \mathbf{h}\, d\Gamma
= \int_{\Gamma_c^{\text{active}}} \mathbf{N}^T
\left(k_p\, g\, \mathbf{n}\right)\, d\Gamma.
\]

In index form, the element residual contribution is
\[
R_{ia}^{(c)}
= \int_{\Gamma_c^{\text{active}}}
N_a\, k_p\, g\, n_i\, d\Gamma.
\]

This is assembled in \texttt{set\_bc\_rigid\_plane\_l(...)} through the local
array \texttt{lR}.

\section{Consistent Linearization Used by the Code}

The code also forms the tensor
\[
\mathbf{n}\otimes\mathbf{n},
\qquad
(\mathbf{n}\otimes\mathbf{n})_{ij} = n_i n_j,
\]
stored locally as \texttt{nDn(i,j)}.

For the standard structural case, the boundary tangent contribution is of the
form
\[
\mathbf{K}_{\Gamma_c}
= \int_{\Gamma_c^{\text{active}}}
k_p\, a_{fu}\, N_a N_b \, (\mathbf{n}\otimes\mathbf{n})\, d\Gamma,
\]
where the generalized-\(\alpha\) factor used in the implementation is
\[
a_{fu} = \texttt{eq.af * eq.beta * dt * dt}.
\]

Thus, the penalty contact adds a rank-one normal-direction stiffness on the
active part of the boundary.

For \texttt{ustruct}, the implementation separates the contribution into
\texttt{lKd} and \texttt{lK}, and the final assembly is passed through
\[
\texttt{ustruct::ustruct\_do\_assem(...)}.
\]
For the standard structural case, the contribution is assembled directly through
\[
\texttt{eq.linear\_algebra->assemble(...)}.
\]

\section{Summary}

In summary:

\begin{itemize}[leftmargin=*]
\item \texttt{RigidPlaneBoundaryCondition.cpp} validates, normalizes, stores,
and distributes the rigid-plane data.
\item The user must provide \texttt{Penalty\_factor}, \texttt{Plane\_point},
\texttt{Plane\_normal}, \texttt{Time\_dependence = Unsteady}, and
\texttt{Temporal\_values\_file\_path}.
\item The actual weak-form modification is assembled in \texttt{set\_bc.cpp},
function \texttt{set\_bc\_rigid\_plane\_l(...)}.
\item The added term is a unilateral penalty traction
\[
\mathbf{h}=k_p g \mathbf{n}
\]
active only when
\[
g=\mathbf{n}\cdot(\mathbf{y}-\mathbf{x}_p)-u_p(t) > 0.
\]
\end{itemize}

\end{document}
